% Generated draft: requires mathematical and source review.
\documentclass[11pt]{article}
\usepackage[margin=25mm]{geometry}
\usepackage{fontspec}
\setmainfont{Latin Modern Roman}
\usepackage{amsmath,amssymb,mathtools,unicode-math}
\setmathfont{Latin Modern Math}
\usepackage{enumitem,needspace,xurl,hyperref}
\setlength{\parindent}{0pt}
\setlength{\parskip}{5pt}
\setlength{\emergencystretch}{3em}
\newcommand{\N}{\mathbb{N}}
\newcommand{\Z}{\mathbb{Z}}
\newcommand{\Q}{\mathbb{Q}}
\newcommand{\R}{\mathbb{R}}
\newcommand{\C}{\mathbb{C}}
\newcommand{\F}{\mathbb{F}}
\newcommand{\A}{\mathbb{A}}
\newcommand{\PP}{\mathbb{P}}
\newcommand{\E}{\mathbb{E}}
\newcommand{\eps}{\varepsilon}
\newcommand{\dd}{\,\mathrm{d}}
\begin{document}
\section*{UnsolvedMath ID: 30005227}
\textbf{Source problem code:} catalogue entry
\subsection{Definitions and mathematical statement}
Let \(\ell\) denote a prime, and let \(\mathbb Z_\ell\) denote the ring of \(\ell\)-adic integers. The supplied problem concerns a spetsial \(\mathbb Z_\ell\)-reflection group \(W\). The source excerpt does not define the spetsial condition or specify the full hypotheses on \(\ell\), \(W\), and the additional data needed for its constructions.

The proposed comparison has two sides. On the representation-theoretic side, the source describes constructing unipotent characters by mimicking Lusztig's theory and defining a principal block by adapting the block theory of Cabanes and Enguehard. On the homotopy-theoretic side, it uses fusion systems constructed by Broto--M\o ller from homotopy fixed points of \(\ell\)-compact groups whose Weyl group is \(W\).

\emph{Source-level problem.} For the spetsial \(\mathbb Z_\ell\)-reflection groups and associated data to which these constructions apply, establish the analogues of the Alperin weight conjecture and Robinson's ordinary weight conjecture that compare invariants of the unipotent principal block with invariants of the associated fusion systems.

\emph{Unresolved formulation.} The supplied excerpt does not give either conjectural identity, define the invariants being compared, specify the homotopy fixed-point maps or their parameters, or state the admissibility hypotheses and quantifiers for these choices. Thus it does not determine a complete mathematical conjecture that can be written as an explicit equality or family of equalities. The precise statements must be recovered from the cited work before a fully quantified formulation is possible; no particular identities or additional hypotheses are asserted here.

The original excerpt reports that these analogues can be formulated and proved in important special cases. It also describes proof ingredients involving generalizations of Yokonuma Hecke algebras to \(\mathbb Z_\ell\)-reflection groups and an analogue of a Morita equivalence due to Puig. These are supplied contextual claims, not additional conjectures or proofs in this draft. The record's dated assessment labels the general problem partially solved, without specifying the established cases or the exact remaining scope.
\subsection{Short English statement}
The problem asks for weight-conjecture analogues connecting unipotent principal-block data associated with a spetsial \(\mathbb Z_\ell\)-reflection group to fusion-system data obtained from homotopy fixed points of an \(\ell\)-compact group with that Weyl group. The supplied excerpt omits the actual identities and their hypotheses, so this draft preserves the research target while recording the information needed for a precise statement.
\subsection{Sources}
\begin{itemize}
\item Local registry: \texttt{data/unsolvedmath/problems.json}, with \texttt{catalog/problems.json} as the supplementary registry; supplied permanent UnsolvedMath ID \(30005227\), source code \texttt{catalogue entry}, title \emph{Weight Conjectures for Spetsial Reflection Groups}. These registry identifiers are supplied provenance; the local files were not inspected for this draft.
\item Supplied report citation: \emph{Character Theory and Categorification} (2023), Oberwolfach Reports, DOI 10.4171/owr/2022/39, pp.~2267--2333. \url{https://doi.org/10.4171/owr/2022/39}. The registry assigns proposed year 2022.
\item Supplied talk attribution: \emph{Weight conjectures for \(\mathbb Z_\ell\)-spetses}, Gunter Malle, joint work with Radha Kessar and Jason Semeraro. The registry does not name a separate proposer.
\item Supplied literature source: Kessar--Malle--Semeraro, weight conjectures. \url{https://doi.org/10.24033/ASENS.2583}.
\item Supplied literature source: Kessar--Malle--Semeraro, preprint. \url{https://arxiv.org/abs/2008.07213}.
\item Supplied literature source: Oberwolfach Report 39/2022. \url{https://ems.press/content/serial-article-files/46975}.
\item Literature assessment dated \texttt{2026-08-22}: the record labels the problem partially solved and reports proofs in significant special cases while describing the general conjectures as unsettled. This is supplied dated curation, not an independently verified literature assessment. The background explicitly requests independent verification of time-sensitive status and cited claims.
\item The registry labels its statement \texttt{corrected\_verified} and reports that the official Oberwolfach report was checked and the question isolated with repaired notation. This is a supplied verification claim; no source visit or independent verification was performed for this draft.
\end{itemize}
\end{document}
